\chapter{Evaluation}
\label{ch:evaluation}

\begin{quote}\itshape
How you find out that yesterday's prompt change made things worse.
\end{quote}

\noindent
Chapter~\ref{ch:observability} ended on a gap. The framework instruments
mechanism thoroughly --- tokens, latency, which executor ran, which tool was
called --- and quality not at all, with the consequence that a green dashboard is
compatible with an agent failing every user. This chapter closes that gap.

It is also the chapter that makes the rest of the book maintainable. Everything
built so far can be changed by editing a string. A prompt is not type-checked, not
covered by a unit test, and not reviewable in any meaningful sense --- you cannot
look at a diff of an instruction and know whether it made the system better. The
only way to know is to measure, and the only way to keep knowing is to measure
automatically.

\begin{verifybox}
Derived from the \code{Evaluation/} folder of \pkg{Microsoft.Agents.AI}, the
workflow evaluation extension in \pkg{Microsoft.Agents.AI.Workflows}, the Foundry
evaluator catalogue, and the evaluation samples upstream. Type and member names
are reliable. The CI configuration in \S\ref{sec:eval-ci} and the trace-mining
code in \S\ref{sec:eval-set} are the author's own and have not been run.
\end{verifybox}

% ============================================================
\section{What ``correct'' means for a non-deterministic system}
\label{sec:eval-correct}
\index{evaluation!determinism}

Start with the thing that makes this hard, because every practice in this chapter
follows from it.

Run the same input through the same agent twice and you may get two different
outputs. Both may be good. One may be good and one may be terrible. So the unit
test contract --- given this input, assert this output --- does not apply, and
teams that try to force it end up with either a test suite asserting on exact
strings that breaks on every model update, or no tests at all. The second is more
common.

What replaces it is a shift in three directions at once.

\textbf{From assertion to scoring.} Not ``is the output equal to this'' but ``how
does this output score against these criteria''. The result is a number or a set
of booleans, not a pass.

\textbf{From one run to many.} A single run tells you almost nothing, because you
cannot tell a systematic change from a sample of one. This is the same discipline
Chapter~\ref{ch:orchestration} \S\ref{sec:orch-benchmark} imposed on the
benchmark, and it is not a coincidence --- \emph{measuring an agent and
benchmarking an agent are the same statistical problem}.

\textbf{From absolute to relative.} ``Is this agent good'' is nearly unanswerable.
``Is this agent better than the one we shipped last week'' is answerable, useful,
and the question you actually have.

\begin{note}
That third shift is the one that makes evaluation practical rather than a
research project. You do not need a scoring system that is correct in an absolute
sense. You need one that is \emph{consistent}, so that a change in the score
means a change in the system. A crude but stable scorer beats a sophisticated
unstable one, every time, and it is available today.
\end{note}

\subsection{Ground truth, where you can get it}

Some things do have a right answer: a classification, an extracted field, a tool
that should have been called. Where ground truth exists, use it --- it is cheaper
and more reliable than anything else in this chapter.

The framework supports this directly. An evaluation item carries an optional
expected output and an optional list of expected tool calls, and reference-based
scorers compare against them. Building a set of a hundred inputs with known
answers is dull work that pays for itself the first time it catches a regression.

The mistake is assuming that because \emph{some} of your system has ground truth,
you can skip evaluating the rest. The parts with no right answer are usually the
parts users notice.

% ============================================================
\section{The evaluation surface}
\label{sec:eval-surface}
\index{evaluation!API}

It is natural to assume that evaluation lives entirely in the extensions
library. It does not, quite. The core agent package owns the abstraction, and
\pkg{Microsoft.Extensions.AI.Evaluation} is one of three things that plug into
it. Knowing that saves you looking in the wrong package for the entry point.

The centre of it is a single interface --- a batch evaluator that takes items and
returns results --- and a set of extension methods that run an agent against
queries and evaluate the results in one call:

\begin{csharp}[caption={Run and evaluate, in one call}, label={lst:eval-basic}]
AgentEvaluationResults results = await agent.EvaluateAsync(queries, evaluator);

Console.WriteLine($"Passed: {results.Passed}/{results.Total}");
\end{csharp}

The same shape works for a completed workflow run, with per-agent breakdown,
which \S\ref{sec:eval-orchestration} uses.

Three implementations of that interface exist, and they form a cost ladder. Pick
deliberately --- the whole art of this chapter is spending model calls only where
cheaper checks cannot reach.

\begin{center}
\small
\begin{tabularx}{\linewidth}{@{}lX@{}}
\toprule
\textbf{Tier} & \textbf{What it is} \\
\midrule
Local checks & Deterministic functions over the response. No model, no network \\[4pt]
Judge model & An extensions-library evaluator, scored by a model you configure \\[4pt]
Foundry & Hosted evaluators run server-side, with a report \\
\bottomrule
\end{tabularx}
\end{center}

\subsection{Local checks: the tier to start with}

A check is a delegate from an evaluation item to a result. That is the whole
abstraction, and it means a custom check is a lambda:

\begin{csharp}[caption={Custom checks, no model required}, label={lst:eval-local}]
EvalCheck noRefusal = FunctionEvaluator.Create("no_refusal", (string response) =>
    !response.Contains("I can't help", StringComparison.OrdinalIgnoreCase)
    && !response.Contains("outside my scope", StringComparison.OrdinalIgnoreCase));

EvalCheck hasSteps = FunctionEvaluator.Create("has_actionable_steps", (string response) =>
    response.Contains("1.", StringComparison.Ordinal)
    || response.Contains("- ", StringComparison.Ordinal));

LocalEvaluator evaluator = new(noRefusal, hasSteps, EvalChecks.NonEmpty(50));

AgentEvaluationResults results = await agent.EvaluateAsync(queries, evaluator);
\end{csharp}

Note what the first of those is. It is the refusal detector
Chapter~\ref{ch:observability} \S\ref{sec:obs-scoring} recommended writing by
hand --- and here it is a two-line lambda, with the framework doing the running
and aggregation. Several more come in the box:

\begin{center}
\small
\begin{tabularx}{\linewidth}{@{}lX@{}}
\toprule
\textbf{Check} & \textbf{Asserts} \\
\midrule
\code{KeywordCheck} & Required keywords appear in the response \\
\code{NonEmpty} & The response is at least a given length \\
\code{ContainsExpected} & The response contains the expected output \\
\code{ToolCalledCheck} & Named tools were called \\
\code{ToolCallsPresent} & Some tool was called at all \\
\code{ToolCallArgsMatch} & Tool arguments match what was expected \\
\code{HasImageContent} & The response includes image content \\
\bottomrule
\end{tabularx}
\end{center}

\begin{note}
\textbf{The tool checks are the underrated ones.} A large share of real agent
failures are not bad prose --- they are the agent not calling the tool it should
have, calling the wrong one, or calling the right one with wrong arguments. All
three are detectable without a model, deterministically, in milliseconds. Build
this tier before you consider a judge.
\end{note}

\subsection{The judge-model tier}

When a check cannot be written as a function --- is this answer relevant, is it
grounded in the context it was given --- the extensions library supplies
evaluators that ask a model. They are adapted onto the same interface, so they
run through the same call, and they need a chat configuration naming the judge
model.

\begin{warning}
A judge model is a language model, with all the properties this chapter opened
by describing. It is non-deterministic, it can be wrong, and it has its own
failure modes --- notably a tendency to score verbose answers higher than terse
correct ones.

Two disciplines make it usable. Use a \emph{different} model from the one under
test, or you are asking a system to grade its own homework. And validate the
judge against human judgement on a sample before you trust it, then re-validate
when you change judge models --- a judge upgrade is a measurement change that
will move every score you have, and if you have not noticed you will read it as a
quality change.
\end{warning}

\subsection{Repetitions are a first-class parameter}

This is the detail that shows the API was designed by someone who had done this
before. The evaluation call takes a repetition count, and running each query $n$
times independently is a supported mode rather than something you build around
the API.

\begin{csharp}[caption={Measuring consistency, not just quality}, label={lst:eval-reps}]
AgentEvaluationResults results = await agent.EvaluateAsync(
    queries, evaluator, numRepetitions: 5);
\end{csharp}

Use it. A query that passes five times out of five and a query that passes three
times out of five are very different states of health, and a single run cannot
distinguish them. This is exactly the discipline \S\ref{sec:orch-benchmark}
demanded of the benchmark, and it is the raw material for
\S\ref{sec:eval-ci}'s answer to flaky gates.

% ============================================================
\section{Foundry evaluations}
\label{sec:eval-foundry}
\index{evaluation!Foundry}

The third tier runs server-side. It implements the same interface, so switching
to it is a change of evaluator and nothing else:

\begin{csharp}[caption={Hosted evaluation}, label={lst:eval-foundry}]
FoundryEvals evaluator = new(projectClient, deploymentName,
    FoundryEvals.Relevance, FoundryEvals.Coherence);

AgentEvaluationResults results = await agent.EvaluateAsync(queries, evaluator);

Console.WriteLine($"Report: {results.ReportUrl}");
\end{csharp}

The catalogue is the reason to care, because it goes well beyond the usual
quality metrics. Grouped by what they measure:

\begin{center}
\small
\begin{tabularx}{\linewidth}{@{}lX@{}}
\toprule
\textbf{Group} & \textbf{Evaluators} \\
\midrule
Agent behaviour & intent resolution, task adherence, task completion, task
navigation efficiency \\[4pt]
Tool use & tool call accuracy, tool selection, tool input accuracy, tool output
utilization, tool call success \\[4pt]
Response quality & coherence, fluency, relevance, groundedness, response
completeness, similarity \\[4pt]
Safety & violence, sexual, self-harm \\
\bottomrule
\end{tabularx}
\end{center}

\begin{note}
Compare that first group against the table in Chapter~\ref{ch:observability}
\S\ref{sec:obs-gap}. ``Whether calling the tool was the right choice'' is tool
selection. ``Whether the turns made progress'' is task navigation efficiency.
``Whether it answered the question asked'' is intent resolution. The things that
chapter listed as uninstrumented are precisely what this catalogue scores --- the
capability is not missing from the ecosystem, it is just not in the telemetry
path, and it costs model calls rather than being free.
\end{note}

Results carry a report URL, a run identifier and a status, because the evaluation
happens as a job rather than in your process. That is the trade: you get a
maintained catalogue and a UI, and you take a network dependency, a per-run cost,
and a requirement that the data leave your process.

\needscreenshot{foundry-evaluation-report}{A hosted evaluation run, scored}{The
Foundry evaluation report page reached from the result's report URL, showing a
completed run with several evaluators scored across a set of queries. At least
one per-item score breakdown must be expanded so the reader can see what an
individual judgement looks like. Browser only.}

\begin{versionbox}
Evaluators are named by string constant, which means the set can grow without a
package update --- and that a name can stop being recognised. Pin the evaluator
names you depend on in configuration, and treat an unrecognised evaluator as a
build failure rather than a silently missing score.
\end{versionbox}

% ============================================================
\section{Building an evaluation set from production traces}
\label{sec:eval-set}
\index{evaluation!building a set}

Everything above needs inputs. Inventing them is the standard approach and it is
the reason most evaluation sets are useless: engineers write the queries they
imagined, which are the ones the system already handles.

Your traces are full of the queries it does not. Chapter~\ref{ch:observability}
put every agent run into a tracing backend; that is an evaluation set waiting to
be extracted.

\subsection{What to select}

Do not sample uniformly. A random sample of production traffic is mostly the easy
case, and evaluating it repeatedly tells you what you already know. Select for
signal:

\begin{itemize}[leftmargin=1.4em]
\item Runs where a cheap scorer from \S\ref{sec:obs-scoring} flagged something ---
      a refusal, a truncation, a tool loop that hit its cap.
\item Runs with unusual shapes: far more turns than the median, far more tokens,
      far longer latency.
\item Runs a user abandoned, if you can tell.
\item Runs that a human later corrected, which are the highest-value items in
      the whole system and are usually thrown away.
\item A modest random sample as a control, so you can tell a regression from a
      change in what you selected.
\end{itemize}

\begin{warning}
\index{evaluation!PII}
Production traces contain user data. An evaluation set built from them is a copy
of that data, living in a repository, being read by a judge model, and quite
possibly leaving your process to a hosted evaluator.

Treat building the set as a data-handling change, not an engineering
convenience: redact before it lands in the repository rather than after,
establish who may read it, and give it a retention policy. This is the second
time this book has raised this --- \S\ref{sec:obs-otel} warned about sensitive
data in spans --- and the two problems compound, because it is the traces with
content enabled that make the best evaluation items.
\end{warning}

\subsection{Curating, which is the actual work}

Extraction gives you inputs. An evaluation set needs inputs \emph{plus} some
notion of what a good response looks like, and that part is human.

The efficient path is to grade a batch by hand once, then use those gradings to
validate a judge, then let the judge scale. What you are buying with the manual
pass is not the labels --- it is the knowledge of whether your scorer agrees with
you.

\begin{note}
Keep the set small enough to run often. Two hundred well-chosen items you run on
every pull request beats ten thousand you run quarterly, because only the first
one changes anyone's behaviour. Grow it by adding cases you got wrong, which is
how a regression suite is supposed to grow anyway.
\end{note}

% ============================================================
\section{Evaluating the orchestration patterns}
\label{sec:eval-orchestration}

Chapter~\ref{ch:orchestration} specified a benchmark and deliberately reported no
results. Its correctness measure was binary --- did the run produce the known
answer --- and \S\ref{sec:orch-benchmark} said plainly that this is not a quality
measure and named this chapter as the place the gap gets closed.

Here is how it closes. A completed workflow run can be evaluated directly, and
the result carries a per-agent breakdown:

\begin{csharp}[caption={Scoring a workflow run, per agent}, label={lst:eval-workflow}]
await using Run run = await InProcessExecution.RunAsync(workflow, task);

AgentEvaluationResults results = await run.EvaluateAsync(evaluator);

Console.WriteLine($"Overall: {results.Passed}/{results.Total}");

foreach (var (agentName, sub) in results.SubResults ?? [])
{
    Console.WriteLine($"  {agentName}: {sub.Passed}/{sub.Total}");
}
\end{csharp}

That per-agent breakdown is what makes the comparison worth running. A pattern's
aggregate score tells you whether it worked; the breakdown tells you \emph{which
participant} was responsible, which is the difference between ``group chat scored
worse'' and ``group chat scored worse because the critic agent contributed
nothing in eight runs out of twenty''. Only the second is actionable.

\begin{note}
\textbf{This is why Chapter~\ref{ch:orchestration} told you to keep the raw event
streams.} Re-scoring a stored run under a new evaluator is cheap. Re-running five
orchestration patterns twenty times each to recover the runs is not --- it is the
entire cost of the benchmark, paid again, because a scorer changed.
\end{note}

\begin{warning}
The tables in Appendix~\ref{app:providers} are still empty. This section
describes how to score the benchmark; it does not report scores, because the
benchmark has not been run. That remains the largest outstanding piece of
original work in this book, and adding an evaluation harness to a study that does
not exist yet does not change its status.
\end{warning}

% ============================================================
\section{Regression gates in CI}
\label{sec:eval-ci}
\index{evaluation!CI gates}

The framework anticipates this, and the evidence is that results carry assertion
methods --- which exist for exactly one reason, which is to be called from a test:

\begin{center}
\small
\begin{tabularx}{\linewidth}{@{}lX@{}}
\toprule
\textbf{Assertion} & \textbf{Fails when} \\
\midrule
\code{AssertAllPassed} & Any item failed any check \\
\code{AssertNoFailedItems} & Any item failed outright \\
\code{AssertScoreAtLeast} & The aggregate score is below a threshold \\
\code{AssertDimensionScoreAtLeast} & One named dimension is below a threshold \\
\bottomrule
\end{tabularx}
\end{center}

\begin{csharp}[caption={An evaluation as a test}, label={lst:eval-test}]
[Fact]
public async Task SupportAgent_MeetsQualityBar()
{
    AgentEvaluationResults results = await this._agent.EvaluateAsync(
        EvaluationSet.Queries, this._evaluator, numRepetitions: 3);

    results.AssertScoreAtLeast(0.8);
}
\end{csharp}

So the mechanism is trivial. The reason most teams do not have this working is
not the mechanism.

\subsection{The flakiness problem, which is the real subject}

Here is what happens. A team adds an evaluation gate. It is non-deterministic, so
it fails occasionally on changes that did nothing wrong. Developers re-run it.
Re-running becomes reflex. Within a month the gate is advisory, and shortly after
that someone removes it because it is slow and always red.

That progression is close to universal, and it is not a discipline failure. It is
the predictable result of putting a stochastic check in a place that demands a
binary answer. Four things actually help.

\textbf{Gate on aggregates, never on individual items.} One item failing is
noise. The mean across two hundred items dropping five points is signal. Assert
on the aggregate and let individual items move.

\textbf{Set the threshold from measured variance, not from a round number.} Run
the suite ten times against an unchanged system. Whatever spread you see is your
noise floor; the gate goes below it. A threshold of 0.8 chosen because it looks
like a good number will fire on noise or miss real regressions, and you will not
know which.

\textbf{Compare against the baseline, not against an absolute.} The question is
whether this change made things worse. Score the merge base, score the branch,
gate on the difference. This also survives model upgrades, which move every
absolute score at once and would otherwise light up every gate you own.

\textbf{Put the expensive tier somewhere other than the pull request.} Local
checks are fast, free and deterministic --- run those on every commit. The judge
tier belongs on a schedule or on merge to the main branch, because a gate that
adds ten minutes and a model bill to every push will be removed, correctly, by
whoever is trying to ship.

\begin{center}
\small
\begin{tabularx}{\linewidth}{@{}lX@{}}
\toprule
\textbf{Where} & \textbf{What runs} \\
\midrule
Every commit & Local checks. Seconds, free, deterministic \\[4pt]
Pull request & Local checks over the full set, plus judge scoring on a small
fixed subset \\[4pt]
Merge to main & Full judge evaluation against the baseline, with repetitions \\[4pt]
Nightly & The whole set, all tiers, trended over time \\
\bottomrule
\end{tabularx}
\end{center}

\begin{note}
The trend line is the thing most teams skip and most regret skipping. A gate
tells you about one change. A chart of the score over sixty days tells you that
quality has been drifting down since a model upgrade three weeks ago --- which no
individual gate would ever have caught, because no single change was responsible.
\end{note}

\needscreenshot{ci-evaluation-gate-failure}{A quality regression, caught before
merge}{A pull request check showing a failed evaluation gate, with the test
output visible: the aggregate score, the threshold it fell below, and at least
two individual item failures listed underneath. The reader must be able to see
that the failure is an aggregate score falling rather than a single assertion.
Browser or terminal, no IDE.}

\begin{exercisebox}
Before Chapter~\ref{ch:hosting}:

\begin{enumerate}[leftmargin=1.4em]
\item Write three local checks for your own agent, one of which is refusal
      detection. Run them over twenty queries. This should take under an hour and
      will find something.
\item Establish your noise floor: run the same suite ten times against an
      unchanged agent and record the spread. Any gate threshold you set without
      this number is guesswork.
\item Take one prompt you believe is an improvement. Score the baseline, score
      the change, and compare. Be honest about the result --- roughly half of
      these turn out to be neutral.
\item Build an evaluation set from real traces using the selection criteria in
      \S\ref{sec:eval-set}. Compare it against a set you wrote from
      imagination. The difference in what fails is the point of the section.
\item Add the cheapest possible gate to CI --- local checks only, aggregate
      threshold, compared against the merge base. Leave it a month and see
      whether anyone has disabled it. If they have, ask why; the answer is
      usually one of the four in \S\ref{sec:eval-ci}.
\end{enumerate}
\end{exercisebox}

% ============================================================
\section{Summary}

Agent output is non-deterministic, so evaluation replaces assertion with scoring,
one run with many, and absolute judgement with comparison against a baseline. A
crude but stable scorer is more useful than a sophisticated unstable one, because
what you need is for a change in the score to mean a change in the system.

Evaluation does not live entirely in the extensions library. The core agent
package owns the abstraction --- a batch evaluator interface, and extension
methods that run an agent or a completed workflow and score the results in one
call --- and there are three implementations forming a cost ladder: local
deterministic checks, judge-model evaluators adapted from the extensions library,
and hosted Foundry evaluators.

Start at the bottom of that ladder. Local checks cost nothing, run in
milliseconds, and the tool-oriented ones catch the failures that matter most ---
the tool not called, the wrong tool called, the right tool called with wrong
arguments. Use a judge only for what a function cannot express, use a different
model from the one under test, and validate the judge against human judgement
before trusting it.

The Foundry catalogue scores precisely what Chapter~\ref{ch:observability} listed
as uninstrumented: intent resolution, task adherence, task navigation efficiency,
tool selection. That capability is not missing from the ecosystem --- it is
outside the telemetry path and costs model calls.

Build evaluation sets from production traces rather than imagination, selecting
for signal rather than sampling uniformly, and treat the set as a data-handling
change because that is what it is.

Repetitions are a first-class parameter and the reason to use them is the same
reason Chapter~\ref{ch:orchestration} demanded distributions rather than best
runs. Workflow runs can be scored with a per-agent breakdown, which turns ``this
pattern scored worse'' into ``this participant contributed nothing'' --- and
which is why the raw event streams from that benchmark are worth keeping.

Gates fail for a predictable reason: a stochastic check placed where a binary
answer is demanded. Gate on aggregates, set the threshold from measured variance,
compare against the baseline rather than an absolute, and keep the expensive tier
off the pull request. Then chart the trend, because the drift no single change
caused is the one nothing else will catch.

Chapter~\ref{ch:hosting} takes all of this somewhere it has to keep running.
